% Generated draft: requires mathematical and source review.
\documentclass[11pt]{article}
\usepackage[margin=25mm]{geometry}
\usepackage{fontspec}
\setmainfont{Latin Modern Roman}
\usepackage{amsmath,amssymb,mathtools,unicode-math}
\setmathfont{Latin Modern Math}
\usepackage{enumitem,needspace,xurl,hyperref}
\setlength{\parindent}{0pt}
\setlength{\parskip}{5pt}
\setlength{\emergencystretch}{3em}
\newcommand{\N}{\mathbb{N}}
\newcommand{\Z}{\mathbb{Z}}
\newcommand{\Q}{\mathbb{Q}}
\newcommand{\R}{\mathbb{R}}
\newcommand{\C}{\mathbb{C}}
\newcommand{\F}{\mathbb{F}}
\newcommand{\A}{\mathbb{A}}
\newcommand{\PP}{\mathbb{P}}
\newcommand{\E}{\mathbb{E}}
\newcommand{\eps}{\varepsilon}
\newcommand{\dd}{\,\mathrm{d}}
\begin{document}
\section*{UnsolvedMath ID: 30006006}
\textbf{Source problem code:} catalogue entry
\subsection{Definitions and mathematical statement}
Let \(p\) be a prime, let \(E/\Q_p\) be a finite extension, let \(G/E\) be a connected reductive quasisplit group, and let \(\ell\ne p\) be a prime. The supplied excerpt describes the following setup for the categorical local Langlands program. Its assertions are retained as source exposition; the record does not supply a separate conjectural statement.

On the automorphic side, write \(\operatorname{Bun}_G\) for the stack of \(G\)-bundles on the Fargues--Fontaine curve. The excerpt says that this stack behaves like a smooth Artin stack of dimension zero and has strata \(\operatorname{Bun}_G^b\) indexed by the Kottwitz set \(B(G)\). For each \(b\in B(G)\), it associates a group \(G_b\), an inner form of a Levi subgroup of \(G\), and describes the corresponding stratum as essentially the classifying stack of the locally profinite group \(G_b(E)\). The word ``essentially'' is a material qualification: the excerpt does not specify an exact identification of stacks.

Using the coefficient notation appearing in the supplied transcription, set
\[
D(\operatorname{Bun}_G):=D_{\mathrm{lis}}(\operatorname{Bun}_G,\Q_\ell).
\]
The excerpt attributes this category of constructible \(\ell\)-adic sheaves to Fargues--Scholze and reports equivalences
\[
D(\operatorname{Bun}_G^b)\simeq D(G_b(E),\Q_\ell)
\qquad (b\in B(G)),
\]
where the right side denotes the derived category of smooth \(\Q_\ell\)-representations of \(G_b(E)\). Writing \(i_b\) for the inclusion of the stratum, the excerpt supplies functors
\[
i_{b!}:D(G_b(E),\Q_\ell)\longrightarrow D(\operatorname{Bun}_G),
\qquad
i_b^*:D(\operatorname{Bun}_G)\longrightarrow D(G_b(E),\Q_\ell).
\]
No adjunction between these two displayed functors is asserted here. The excerpt further reports a semiorthogonal decomposition into the categories \(D(G_b(E),\Q_\ell)\), without specifying its ordering or precise categorical conventions. For the neutral class \(b=1\), it states that \(G_1=G\) and that \(i_{1!}\) embeds smooth \(G(E)\)-representations fully faithfully into sheaves on \(\operatorname{Bun}_G\).

On the spectral side, let \(W_E\) denote the Weil group of \(E\), let \(\widehat G\) denote the dual group, and retain the source notation \({}^{L}G\) for the \(L\)-group. The excerpt describes \(\operatorname{Par}_G\) as a stack of \(\ell\)-adically continuous \(L\)-parameters
\[
\phi:W_E\longrightarrow {}^{L}G(\Q_\ell).
\]
It explicitly acknowledges that making this definition precise is subtle, but does not supply that definition. Attributing the results to independent works of Fargues--Scholze, Zhu, Hellmann, and Dat--Helm--Kurinczuk--Moss, it reports that \(\operatorname{Par}_G\) is a reduced Artin stack, is globally a local complete intersection of pure dimension zero over \(\operatorname{Spec}\Q_\ell\), and has each connected component expressible as a quotient of an affine variety by a reductive group action. It also reports a canonical map
\[
\operatorname{Par}_G\longrightarrow B\widehat G,
\]
where \(B\widehat G\) denotes the classifying stack of \(\widehat G\), distinct from the Kottwitz set \(B(G)\).

\textbf{Unresolved mathematical target.} David Hansen describes joint work in progress with Lucas Mann whose goal is to formulate a program proving the categorical local Langlands conjecture of Fargues--Scholze, cited in the excerpt as reference [3], for many groups. Neither the conjecture nor the intended class of ``many groups'' is specified in this record. The supplied registry marks the statement as \texttt{unrecoverable} and says that the formal conjecture appears instead in source record \texttt{OWR-14298588-008}, whose contents are not supplied here. Consequently, the setup above cannot be promoted to a precise conjecture, a proposed equivalence of categories, or a universally quantified assertion for all such \(G\). The technical definitions, hypotheses, and target assertion remain subject to human mathematical review.
\subsection{Short English statement}
This record preserves David Hansen's introductory account of joint work with Lucas Mann on the categorical local Langlands program. It describes automorphic and spectral objects for a connected reductive quasisplit group over a finite extension of \(\Q_p\), with \(\ell\ne p\). The supplied excerpt does not state the conjecture or specify the groups targeted by the program, so a separate precise problem cannot be recovered. The supplied \texttt{partially\_solved} label is dated curation, not an independently verified mathematical conclusion.
\subsection{Sources}
\begin{itemize}
\item Supplied local UnsolvedMath registry record: permanent ID \(30006006\), title ``Categorical Local Langlands Program Summary,'' exact source code \texttt{catalogue entry}, in the Oberwolfach Reports Open Problems collection. The record lists proposed year 2024, status \texttt{partially\_solved}, and statement status \texttt{unrecoverable}; no clean statement is supplied.
\item Supplied report citation: \emph{Arithmetic Geometry (2025)}, Oberwolfach Reports, DOI 10.4171/owr/2024/33, pp. 1855--1912, \url{https://doi.org/10.4171/owr/2024/33}. The excerpt identifies the talk as ``The categorical local Langlands program,'' by David Hansen, joint work with Lucas Mann, and describes the work as in progress.
\item The excerpt attributes the categorical local Langlands conjecture and the sheaf-category construction to Fargues--Scholze, referring to an unavailable reference [3]. It attributes the stated properties of the parameter stack to independent works of Fargues--Scholze, Zhu, Hellmann, and Dat--Helm--Kurinczuk--Moss. Full bibliographic entries for those works are not supplied.
\item Supplied literature source: Hansen--Mann, ``categorical local Langlands,'' \url{https://arxiv.org/abs/2606.00983}.
\item Supplied literature source: Helm--Solleveld--Xu, ``generic categorical LLC,'' \url{https://arxiv.org/abs/2606.10149}.
\item Literature assessment dated exactly \texttt{2026-08-22}: the supplied curation describes substantial conditional and unconditional results for GL-type, generic, and inductively accessible groups, while stating that the full conjecture for every connected reductive quasisplit \(G\) is not proved. This is supplied dated curation derived from the OWR record, not an independently verified assessment; its time-sensitive status and cited claims require independent verification.
\item The supplied statement-verification note says that the official report identifies this extraction as introductory programmatic exposition and locates the formal conjecture in \texttt{OWR-14298588-008}. This is a supplied verification claim; neither that other record nor the report has been independently consulted for this draft.
\end{itemize}
\end{document}
